\documentclass[10pt,a4paper]{amsart}
\usepackage[margin=19mm]{geometry}
\usepackage{amsmath,amssymb,mathtools}
\usepackage[scr=rsfs,cal=euler]{mathalfa}
\usepackage{xfrac}
\usepackage[unicode,hidelinks,bookmarks=false]{hyperref}
\newcommand{\ie}{i.e.\ }
\newcommand{\cf}{cf.\ }
\newcommand{\resp}{respectively\ }
\newcommand{\Subsection}{\subsection}
\providecommand{\nobreakdash}{\nobreak\hbox{-}}
\setlength{\emergencystretch}{2em}
\multlinegap0pt
\usepackage{bm}
\usepackage[shortlabels]{enumitem}
\usepackage{amssymb}
\usepackage{mathtools}
\usepackage{bbm}
\usepackage{tikz}
\usepackage{stackengine}
\newtheorem{notation}{Notation}
\newtheorem{prop}{Proposition}
\newtheorem{cor}{Corollary}
\newtheorem{definition}{Definition}
\newtheorem{lem}{Lemma}
\usepackage{cleveref}
\crefname{notation}{Notation}{Notations}
\crefname{prop}{Proposition}{Propositions}
\crefname{cor}{Corollary}{Corollarys}
\crefname{definition}{Definition}{Definitions}
\crefname{lem}{Lemma}{Lemmas}
\newcommand{\NN}{\mathbb N}
\newcommand{\PP}{\mathbb P}
\newcommand{\bfx}{\boldsymbol{x}}
\title{A probabilistic interpretation for\\ interpolation Macdonald polynomials}
\author{Houcine Ben Dali, Lauren Williams}
\date{Source: arXiv:2602.13492v1 (CC BY 4.0)}
\begin{document}
\maketitle

\noindent\emph{Source and licence.} A probabilistic interpretation for\\ interpolation Macdonald polynomials. Version arXiv:2602.13492v1 (CC BY 4.0), distributed by arXiv under the Creative Commons Attribution 4.0 International licence, http://creativecommons.org/licenses/by/4.0/, as recorded in the arXiv licence metadata for this exact version and shown on the abstract page https://arxiv.org/abs/2602.13492v1. Full licence notice: \texttt{licenses/arxiv-2602.13492-CC-BY-4.0.txt}.\par\smallskip

\noindent\textit{Editorial scope.} Counted unit: the article's prop prop:density (source0.tex lines 1144--1157). The article's complete original proof of it is delivered with it (source0.tex lines 1159--1177, one \texttt{proof} environment of 19 lines). No mathematics is written, repaired or re-derived: the statement and the proof are the article's own lines, in the article's own notation, and no retained line is substituted or re-flowed.  Retained uncounted as the source-local context the proof needs: lines 301--312; lines 745--750; lines 789--794; lines 1117--1120; lines 1129--1132; lines 1136--1142.  Nothing was omitted from any retained span, so the package carries no omission record.  No retained line contains a citation, so the wrapper carries no bibliography and no bibliography entry.  No retained line contains a figure environment, an image-inclusion command or drawing code, so no image is delivered and the package declares no asset dependency.  Editorial formatting only, in addition: an AMS \texttt{amsart} preamble that restates the article's own declarations and macros used by the retained lines, namely the environments \texttt{notation}, \texttt{prop}, \texttt{cor}, \texttt{definition}, \texttt{lem} on separate counters, together with the standard packages those lines need.  The retained environments print in this document as Notation 1, Proposition 1, Corollary 1, Definition 1, Lemma 1 in order of appearance, whereas the article prints the same items with the numbering of its own full document.  The complete native source of the article as distributed in the arXiv e-print for this exact version is bundled unchanged as \texttt{sources/194644/source0.tex}.
\begin{notation}\label{not:type}
Let $\lambda=(\lambda_1,\dots, \lambda_n)$
with $\lambda_1\geq\dots\geq\lambda_n\geq 0$ be a partition. We can
describe such a partition by its vector of types
$\mathbf{m}=(m_0, m_1, \dots, m_L)$, where
$m_i=m_i(\lambda):=\#\{j:\lambda_j=i\}$,
and $L$ is the largest part that occurs.
We thus sometimes denote our partition by
$\lambda=(L^{m_L}, \dots, 1^{m_1}, 0^{m_0})$.
Note that $\sum_{i=0}^L m_i=n$.  We also write
$M_i:=m_i+ m_{i+1} + \dots + m_L$.
\end{notation}

\begin{prop}\label{prop:lumping}
    Let $\phi:\NN \to \NN$ be a weakly order-preserving function with $\phi(0)=0$, and let $\lambda=(\lambda_1,\dots,\lambda_n)$ and $\kappa=(\kappa_1,\dots,\kappa_n)$ be  partitions such that $\kappa=\phi(\lambda)$.
Then the interpolation $t$-Push TASEP with particle content $\lambda$ lumps (or projects) to the interpolation $t$-Push TASEP with particle content $\kappa$.  That is,  
    if $\eta,\zeta\in S_n(\kappa)$ and $\mu\in S_n(\lambda)$ such that $\phi(\mu)=\eta$, then
    $$\sum_{\nu\in S_n(\lambda):\, \phi(\nu)=\zeta}\PP_{\lambda}(\mu,\nu)=\PP_\kappa(\eta,\zeta).$$
\end{prop}

\begin{cor}\label{cor:lumping}
Let $\phi,\lambda$ and $\kappa$ as above. Then  the stationary distributions of the 
interpolation $t$-Push TASEPs with particle contents $\kappa$ and $\lambda$ are related as follows:
for any $\eta\in S_n(\kappa)$, 
$$\pi^*_\kappa(\eta)=\sum_{\rho\in S_n(\lambda):\, \phi(\rho)=\eta}\pi^*_\lambda(\rho).$$
\end{cor}

  \begin{equation}\label{eq:density}
  \langle \eta_1 \rangle =
  \frac{(x_1-\frac{1}{t^{n-1}}) e^*_{m_1-1}(tx_2, tx_3,\dots,tx_n;t)} {t^{m_1-1} e^*_{m_1}(x_1,\dots,x_n;t)},
  \end{equation}

\begin{definition}\label{def:sstar}
For any partition $\lambda=(\lambda_1,\dots,\lambda_n)$, we define the \emph{$t$-interpolation Schur polynomial $s_\lambda^*(x_1,\dots,x_n;t)=s^*_\lambda(\bfx;t)$} by
    $$s^*_\lambda(\bfx;t)=P^*_\lambda(\bfx;t,t).$$
\end{definition}

\begin{lem}\label{lem:2column}
Let $\lambda=(2^a,1^b)$.  Then we have 
\begin{align*}
s^*_{\lambda}(t\bfx;t)&=e^*_{a+b,n-1}(t\bfx;t)e^*_{a,n-2}(t\bfx;t)-e^*_{a+b+1,n-2}(t\bfx;t)e^*_{a-1,n-1}(t\bfx;t)\\
&=t^a e^*_{a+b}(t\bfx;t) e^*_a(\bfx;t) - t^{a+b+1} e^*_{a+b+1}(\bfx;t) e^*_{a-1}(t\bfx;t).
\end{align*}
\end{lem}

\begin{prop}\label{prop:density}
Recall, as in \cref{not:type}, that  $M_i = m_i+m_{i+1} + \dots + m_L.$  Let $s^*_{\lambda}$ be the $t$-interpolation Schur polynomial from \cref{def:sstar}.
   The density of the particle of species $i$ in the first site in 
   the interpolation $t$-Push TASEP with content $\lambda=(L^{m_L}, \dots, 1^{m_1}, 0^{m_0})$ is given by
   \small{
   \begin{align*}
       \langle \eta_1^{(i)} \rangle =&
       \frac{(x_1-\frac{1}{t^{n-1}}) \left(e_{M_i-1}^*\left(t\bfx^{(1)};t\right) e^*_{M_{i+1}}\left(\bfx^{(1)};t\right)  - t^{m_i} e_{M_{i+1}-1}^*\left(t\bfx^{(1)};t\right) e_{M_i}^*\left(\bfx^{(1)};t\right) \right)}{t^{M_{i}-1} e^*_{M_i}\left(\bfx;t\right) e^*_{M_{i+1}}\left(\bfx;t\right)} \\
       =&
       \frac{(x_1-\frac{1}{t^{n-1}}) 
       \ s^*_{(2^{M_{i+1}},1^{m_i-1})}\left(t\bfx^{(1)};t\right)
       }{t^{M_{i+1}+M_i-1} e^*_{M_i}\left(\bfx;t\right)e^*_{M_{i+1}}\left(\bfx;t\right)}.
   \end{align*}}
\end{prop}

\begin{proof}
By \cref{prop:lumping} and the lumping property (\cref{cor:lumping}),  the density of the particle
    of species $i$ is the density of the particle of species $1$ in the single species
    Markov chain with $M_i$ particles, minus the density of the particle of species $1$
    in the single species Markov chain with $M_{i+1}$ particles.
    Thus using \cref{eq:density}, we have that 
     \begin{align*}
       \langle \eta_1^{(i)} \rangle &= 
       \frac{(x_1-\frac{1}{t^{n-1}}) e^*_{M_i-1}\left(t\bfx^{(1)};t\right)}{t^{M_i-1} e^*_{M_i}(\bfx;t)}  - \frac{(x_1-\frac{1}{t^{n-1}}) e^*_{M_{i+1}-1}\left(t\bfx^{(1)};t\right)}{t^{M_{i+1}-1} e^*_{M_{i+1}}(\bfx;t)} \\
       &= \frac{(x_1-\frac{1}{t^{n-1}})}{t^{M_{i}-1}} \cdot 
       \frac{e^*_{M_i-1}\left(t\bfx^{(1)};t\right) e^*_{M_{i+1}}(\bfx;t)-t^{m_i} e^*_{M_{i+1}-1}\left(t\bfx^{(1)};t\right) e^*_{M_i}(\bfx;t)}{e^*_{M_i}(\bfx;t) e^*_{M_{i+1}}(\bfx;t)}, \label{eq:sofar}
   \end{align*}
   where $\bfx:=(x_1,\dots,x_n)$.
If we now use the fact that 
\begin{equation}
e^*_d(\bfx;t) = \left(x_1-\frac{1}{t^{n-1}}\right) t^{-(d-1)} e^*_{d-1}\left(t\bfx^{(1)};t\right) + e^*_d\left(\bfx^{(1)};t\right),    
\end{equation}
we get the first statement of the proposition.  To get the second statement, we use \cref{lem:2column}.
\end{proof}

\end{document}
